


Durante esta tarea se estudia experimentalmente cómo se comportan distintos algoritmos para dos problemas específicos: El ordenamiento de arreglos en una dimensión y la multiplicación de matrices cuadradas entre ellas. Para el primer problema se utilizaron y compararon los siguietnes algoritmos: MergeSort, QuickSort, PatienceSort y sort, mientras que para la multiplicación de matrices se utiliza Naive y Strassen. su objetivo es analizar los tiempos de ejecucion y tambien de uso de memoria para diferentes tamaños y/o caracteristicas en las entradas, asi poder contrastar los resultados, la velocidad de obtencion y el espacio necesario para ser ejecutado, todo esto dependiendo de la complejidad de cada algoritmo.

